% ---------------------------------------------------------------------
%  5. Resultados — Integrante C
%  Todos los valores provienen de results/ (CSV y macros generados por
%  BenchmarkRunner); build.sh copia ../results en report/results.
% ---------------------------------------------------------------------
\section{Resultados}
\label{sec:resultados}

% Gráfico de E1: tiempo frente a n en escala log-log.
% #1 = archivo CSV, #2 = título, #3 = extremos de la recta de pendiente 1.
\newcommand{\graficoEuno}[3]{%
  \begin{tikzpicture}
    \begin{loglogaxis}[informe, height=0.78\linewidth, log basis x=2,
        xlabel={$n$ [bytes]}, ylabel={Mediana [ms]},
        title={#2}, title style={font=\small\bfseries},
        legend to name=leyendaEuno, legend columns=5,
        legend style={draw=gray!50, /tikz/every even column/.append style={column sep=6pt}}]
      \addplot+[huffenc, error bars/.cd, y dir=plus, y explicit]
        table[col sep=comma, x=n, y=huff_enc_ms, y error=huff_enc_ci] {results/#1};
      \addlegendentry{Huffman codificación}
      \addplot+[huffdec, error bars/.cd, y dir=plus, y explicit]
        table[col sep=comma, x=n, y=huff_dec_ms, y error=huff_dec_ci] {results/#1};
      \addlegendentry{Huffman decodificación}
      \addplot+[lzwenc, error bars/.cd, y dir=plus, y explicit]
        table[col sep=comma, x=n, y=lzw_enc_ms, y error=lzw_enc_ci] {results/#1};
      \addlegendentry{LZW codificación}
      \addplot+[lzwdec, error bars/.cd, y dir=plus, y explicit]
        table[col sep=comma, x=n, y=lzw_dec_ms, y error=lzw_dec_ci] {results/#1};
      \addlegendentry{LZW decodificación}
      \addplot[pendiente] coordinates {#3};
      \addlegendentry{pendiente 1 (referencia)}
    \end{loglogaxis}
  \end{tikzpicture}}

% Gráfico de E1: costo por byte frente a n.
\newcommand{\graficoNspb}[2]{%
  \begin{tikzpicture}
    \begin{semilogxaxis}[informe, height=0.78\linewidth, log basis x=2,
        xlabel={$n$ [bytes]}, ylabel={ns/byte}, ymin=0,
        title={#2}, title style={font=\small\bfseries}, legend pos=north east]
      \addplot[huffenc] table[col sep=comma, x=n, y=huff_enc_nspb] {results/#1};
      \addlegendentry{Huffman cod.}
      \addplot[huffdec] table[col sep=comma, x=n, y=huff_dec_nspb] {results/#1};
      \addlegendentry{Huffman decod.}
      \addplot[lzwenc] table[col sep=comma, x=n, y=lzw_enc_nspb] {results/#1};
      \addlegendentry{LZW cod.}
      \addplot[lzwdec] table[col sep=comma, x=n, y=lzw_dec_nspb] {results/#1};
      \addlegendentry{LZW decod.}
    \end{semilogxaxis}
  \end{tikzpicture}}

% Gráfico de E1: bits por byte frente a n para un algoritmo (#1 = prefijo de columna).
\newcommand{\graficoBpbN}[2]{%
  \begin{tikzpicture}
    \begin{semilogxaxis}[informe, height=0.78\linewidth, log basis x=2,
        xlabel={$n$ [bytes]}, ylabel={bits por byte}, ymin=0,
        title={#2}, title style={font=\small\bfseries},
        legend to name=leyendaBpbN, legend columns=4,
        legend style={draw=gray!50, /tikz/every even column/.append style={column sep=6pt}}]
      \addplot[black!70, mark=*, mark size=1.2pt] table[col sep=comma, x=n, y=#1_bpb] {results/e1_uniform.csv};
      \addlegendentry{\texttt{uniform}}
      \addplot[blue!70!black, mark=square*, mark size=1.2pt] table[col sep=comma, x=n, y=#1_bpb] {results/e1_skewed.csv};
      \addlegendentry{\texttt{skewed}}
      \addplot[green!45!black, mark=triangle*, mark size=1.5pt] table[col sep=comma, x=n, y=#1_bpb] {results/e1_synthtext.csv};
      \addlegendentry{\texttt{synthtext}}
      \addplot[red!70!black, mark=diamond*, mark size=1.5pt] table[col sep=comma, x=n, y=#1_bpb] {results/e1_dickens.csv};
      \addlegendentry{\texttt{dickens}}
    \end{semilogxaxis}
  \end{tikzpicture}}

\subsection{E1: escalamiento del tiempo con el tamaño de la entrada}
\label{sec:res-e1}
\begin{aporte}{C}
La \cref{fig:e1-tiempo} muestra la mediana del tiempo de codificación y de decodificación en
función de $n$ para las cuatro fuentes, en escala log--log; la recta discontinua tiene
pendiente 1 y sirve solo como guía visual de crecimiento lineal (su posición vertical es
arbitraria). Las barras de error indican el semiancho superior del intervalo de confianza del
95\,\%.

\siresultado{e1_dickens.csv}{%
\begin{figure}[H]
  \centering
  \begin{subfigure}[t]{0.49\linewidth}
    \centering
    \graficoEuno{e1_uniform.csv}{\texttt{uniform}}{(16384,0.16384) (16777216,167.77216)}
    \caption{Bytes uniformes i.i.d.}
    \label{fig:e1-uniform}
  \end{subfigure}\hfill
  \begin{subfigure}[t]{0.49\linewidth}
    \centering
    \graficoEuno{e1_skewed.csv}{\texttt{skewed}}{(16384,0.16384) (16777216,167.77216)}
    \caption{Geométrica i.i.d. ($p = 0.15$).}
    \label{fig:e1-skewed}
  \end{subfigure}\\[6pt]
  \begin{subfigure}[t]{0.49\linewidth}
    \centering
    \graficoEuno{e1_synthtext.csv}{\texttt{synthtext}}{(16384,0.16384) (16777216,167.77216)}
    \caption{Texto sintético (Zipf).}
    \label{fig:e1-synthtext}
  \end{subfigure}\hfill
  \begin{subfigure}[t]{0.49\linewidth}
    \centering
    \graficoEuno{e1_dickens.csv}{\texttt{dickens}}{(16384,0.16384) (8388608,83.88608)}
    \caption{Texto real (prefijos de \texttt{dickens}).}
    \label{fig:e1-dickens}
  \end{subfigure}\\[4pt]
  \pgfplotslegendfromname{leyendaEuno}
  \caption{E1: mediana del tiempo frente al tamaño de la entrada (escala log--log).}
  \label{fig:e1-tiempo}
\end{figure}}

% TODO-RESULTS: describir la forma de las curvas de la figura e1-tiempo (paralelismo con la recta de pendiente 1, desviaciones en tamaños pequeños, diferencias entre fuentes y entre codificación/decodificación).
\todoresults{Interpretar la \cref{fig:e1-tiempo}: paralelismo de las curvas con la recta de
pendiente 1, desviaciones en los tamaños pequeños y diferencias entre fuentes y entre
codificación y decodificación.}

La \cref{tab:pendientes} reporta la pendiente del ajuste log--log y su $R^2$ para cada fuente,
algoritmo y operación, calculados con $n \ge 2^{16}$, junto con el costo por byte en el mayor
tamaño medido.

\siresultado{fits.csv}{%
\begin{table}[H]
  \centering
  \caption{E1: pendiente $\hat{k}$ del ajuste $\ln t = \hat{k}\ln n + c$, coeficiente $R^2$ y
           costo por byte en el mayor $n$.}
  \label{tab:pendientes}
  \small
  \pgfplotstabletypeset[
    col sep=comma,
    columns={source,algorithm,operation,slope,r2,nspb_at_max},
    columns/source/.style={string type, column name={Fuente},
      postproc cell content/.append style={/pgfplots/table/@cell content/.add={\ttfamily}{}}},
    columns/algorithm/.style={string type, column name={Algoritmo},
      string replace={huffman}{Huffman}, string replace={lzw}{LZW}},
    columns/operation/.style={string type, column name={Operación},
      string replace={encode}{codificación}, string replace={decode}{decodificación}},
    columns/slope/.style={column name={Pendiente $\hat{k}$}, fixed, fixed zerofill, precision=3},
    columns/r2/.style={column name={$R^2$}, fixed, fixed zerofill, precision=4},
    columns/nspb_at_max/.style={column name={ns/byte en $n_{\max}$}, fixed, fixed zerofill, precision=2},
    every head row/.style={before row=\toprule, after row=\midrule},
    every last row/.style={after row=\bottomrule},
    every row no 3/.style={after row=\midrule},
    every row no 7/.style={after row=\midrule},
    every row no 11/.style={after row=\midrule},
  ]{results/fits.csv}
\end{table}}

% TODO-RESULTS: comentar el rango de pendientes y R^2 de la tabla tab:pendientes (por ejemplo, con \resSlopeDickensHuffEnc, \resSlopeDickensLzwEnc) y su cercanía a 1.
\todoresults{Comentar el rango de pendientes y de $R^2$ de la \cref{tab:pendientes} y su
cercanía a 1 (por ejemplo, en \texttt{dickens}: Huffman \resSlopeDickensHuffEnc{} /
\resSlopeDickensHuffDec{}, LZW \resSlopeDickensLzwEnc{} / \resSlopeDickensLzwDec{}).}

La \cref{fig:e1-nspb} presenta el mismo experimento como costo por byte en escala
semilogarítmica: bajo complejidad lineal, las curvas deben ser aproximadamente horizontales.

\siresultado{e1_synthtext.csv}{%
\begin{figure}[H]
  \centering
  \begin{subfigure}[t]{0.49\linewidth}
    \centering
    \graficoNspb{e1_dickens.csv}{\texttt{dickens}}
    \caption{Texto real.}
    \label{fig:e1-nspb-dickens}
  \end{subfigure}\hfill
  \begin{subfigure}[t]{0.49\linewidth}
    \centering
    \graficoNspb{e1_synthtext.csv}{\texttt{synthtext}}
    \caption{Texto sintético.}
    \label{fig:e1-nspb-synthtext}
  \end{subfigure}
  \caption{E1: costo por byte (ns/byte) frente al tamaño de la entrada.}
  \label{fig:e1-nspb}
\end{figure}}

% TODO-RESULTS: interpretar la figura e1-nspb (estabilidad del costo por byte, efectos de caché o de costos fijos en tamaños pequeños, comparación del costo relativo de ambos algoritmos).
\todoresults{Interpretar la \cref{fig:e1-nspb}: estabilidad del costo por byte, efectos de
costos fijos o de caché y costo relativo de cada algoritmo y operación.}

Finalmente, la \cref{fig:e1-bpb} muestra la tasa de compresión obtenida en E1 para cada
fuente y tamaño.

\siresultado{e1_uniform.csv}{%
\begin{figure}[H]
  \centering
  \begin{subfigure}[t]{0.49\linewidth}
    \centering
    \graficoBpbN{huff}{Huffman}
    \caption{Huffman.}
    \label{fig:e1-bpb-huff}
  \end{subfigure}\hfill
  \begin{subfigure}[t]{0.49\linewidth}
    \centering
    \graficoBpbN{lzw}{LZW ($b = 16$)}
    \caption{LZW.}
    \label{fig:e1-bpb-lzw}
  \end{subfigure}\\[4pt]
  \pgfplotslegendfromname{leyendaBpbN}
  \caption{E1: bits por byte frente al tamaño de la entrada.}
  \label{fig:e1-bpb}
\end{figure}}

% TODO-RESULTS: interpretar la figura e1-bpb (efecto de la cabecera de Huffman en n pequeño, aprendizaje del diccionario de LZW, expansión de LZW en uniform).
\todoresults{Interpretar la \cref{fig:e1-bpb}: efecto de la cabecera de 262 bytes de Huffman en
$n$ pequeño, mejora de LZW a medida que el diccionario aprende y comportamiento en
\texttt{uniform}.}
\end{aporte}

\subsection{E2: barrido de entropía}
\label{sec:res-e2}
\begin{aporte}{C}
En E2 la fuente es i.i.d. uniforme sobre $k$ símbolos, de modo que $H_0 = \log_2 k$ se conoce
por construcción y no hay frases repetidas más allá de las debidas al azar. La
\cref{fig:e2-bpb} compara los bits por byte de ambos algoritmos con la entropía empírica y con
la cota superior $H_0 + 1$ de la \cref{eq:cota-huffman}.

\siresultado{e2_entropy.csv}{%
\begin{figure}[H]
  \centering
  \begin{tikzpicture}
    \begin{semilogxaxis}[informe, width=0.72\linewidth, height=0.5\linewidth, log basis x=2,
        xlabel={$k$ (tamaño del alfabeto); $H_0 = \log_2 k$}, ylabel={bits por byte},
        ymin=0, legend pos=north west]
      \addplot[huffenc] table[col sep=comma, x=k, y=huff_bpb] {results/e2_entropy.csv};
      \addlegendentry{Huffman}
      \addplot[lzwenc] table[col sep=comma, x=k, y=lzw_bpb] {results/e2_entropy.csv};
      \addlegendentry{LZW ($b = 16$)}
      \addplot[black, thick, mark=none] table[col sep=comma, x=k, y=h0] {results/e2_entropy.csv};
      \addlegendentry{$H_0$ empírica}
      \addplot[pendiente] table[col sep=comma, x=k, y expr=\thisrow{h0}+1] {results/e2_entropy.csv};
      \addlegendentry{$H_0 + 1$}
    \end{semilogxaxis}
  \end{tikzpicture}
  \caption{E2: bits por byte frente al tamaño del alfabeto ($n = 2^{23}$).}
  \label{fig:e2-bpb}
\end{figure}}

% TODO-RESULTS: interpretar e2-bpb (Huffman frente a [H0, H0+1), distancia de LZW a H0 por k; citar por ejemplo \resEntropyHuffBpbKTwo, \resEntropyLzwBpbKTwo, \resEntropyHuffBpbKTwoFiftySix, \resEntropyLzwBpbKTwoFiftySix).
\todoresults{Interpretar la \cref{fig:e2-bpb}: posición de Huffman en $[H_0, H_0+1)$ y
distancia de LZW a $H_0$ para cada $k$ (por ejemplo, $k = 2$: Huffman \resEntropyHuffBpbKTwo{},
LZW \resEntropyLzwBpbKTwo{}; $k = 256$: Huffman \resEntropyHuffBpbKTwoFiftySix{}, LZW
\resEntropyLzwBpbKTwoFiftySix{}).}

La \cref{fig:e2-tiempo} muestra cómo varía el tiempo con la entropía para un tamaño fijo, lo
que permite contrastar el término $\Lbar$ del costo de decodificación de Huffman.

\siresultado{e2_entropy.csv}{%
\begin{figure}[H]
  \centering
  \begin{subfigure}[t]{0.49\linewidth}
    \centering
    \begin{tikzpicture}
      \begin{semilogxaxis}[informe, height=0.78\linewidth, log basis x=2,
          xlabel={$k$}, ylabel={Mediana [ms]}, ymin=0, legend pos=north west]
        \addplot+[huffdec, error bars/.cd, y dir=plus, y explicit]
          table[col sep=comma, x=k, y=huff_dec_ms, y error=huff_dec_ci] {results/e2_entropy.csv};
        \addlegendentry{Huffman}
        \addplot+[lzwdec, error bars/.cd, y dir=plus, y explicit]
          table[col sep=comma, x=k, y=lzw_dec_ms, y error=lzw_dec_ci] {results/e2_entropy.csv};
        \addlegendentry{LZW}
      \end{semilogxaxis}
    \end{tikzpicture}
    \caption{Decodificación.}
    \label{fig:e2-dec}
  \end{subfigure}\hfill
  \begin{subfigure}[t]{0.49\linewidth}
    \centering
    \begin{tikzpicture}
      \begin{semilogxaxis}[informe, height=0.78\linewidth, log basis x=2,
          xlabel={$k$}, ylabel={Mediana [ms]}, ymin=0, legend pos=north west]
        \addplot+[huffenc, error bars/.cd, y dir=plus, y explicit]
          table[col sep=comma, x=k, y=huff_enc_ms, y error=huff_enc_ci] {results/e2_entropy.csv};
        \addlegendentry{Huffman}
        \addplot+[lzwenc, error bars/.cd, y dir=plus, y explicit]
          table[col sep=comma, x=k, y=lzw_enc_ms, y error=lzw_enc_ci] {results/e2_entropy.csv};
        \addlegendentry{LZW}
      \end{semilogxaxis}
    \end{tikzpicture}
    \caption{Codificación.}
    \label{fig:e2-enc}
  \end{subfigure}
  \caption{E2: tiempo frente al tamaño del alfabeto ($n = 2^{23}$).}
  \label{fig:e2-tiempo}
\end{figure}}

% TODO-RESULTS: interpretar e2-tiempo (crecimiento del tiempo de decodificación de Huffman con log2 k frente a LZW; citar \resEntropyHuffDecMsKTwo, \resEntropyHuffDecMsKTwoFiftySix, \resEntropyLzwDecMsKTwo, \resEntropyLzwDecMsKTwoFiftySix).
\todoresults{Interpretar la \cref{fig:e2-tiempo}: tendencia del tiempo de decodificación de
Huffman con $\log_2 k$ frente a LZW (por ejemplo, Huffman \resEntropyHuffDecMsKTwo{} ms con
$k = 2$ y \resEntropyHuffDecMsKTwoFiftySix{} ms con $k = 256$; LZW \resEntropyLzwDecMsKTwo{} y
\resEntropyLzwDecMsKTwoFiftySix{} ms).}
\end{aporte}

\subsection{E3: corpus Canterbury y Silesia}
\label{sec:res-e3}
\begin{aporte}{C}
La \cref{tab:corpus} reporta, para cada archivo, su tamaño, su entropía de orden cero, los bits
por byte de cada algoritmo y el rendimiento de codificación. En negrita aparece el mejor de los
dos algoritmos comparados; DEFLATE se muestra solo como referencia.

\siresultado{table_corpus.tex}{%
\begin{table}[H]
  \centering
  \caption{E3: bits por byte y rendimiento de codificación en los corpus Canterbury (arriba) y
           Silesia (abajo).}
  \label{tab:corpus}
  \small
  \setlength{\tabcolsep}{5pt}
  \begin{tabular}{@{}l r r r r r r r@{}}
    \toprule
    Archivo & Tamaño (KiB) & $H_0$ & Huffman & LZW & DEFLATE$^\dagger$ & Huff. MB/s & LZW MB/s \\
    \midrule
    \texttt{E.coli} & 4530 & 2.00 & \textbf{2.00} & 2.17 & 2.31 & 348 & 141 \\
\texttt{alice29.txt} & 149 & 4.57 & 4.63 & \textbf{3.27} & 2.86 & 222 & 105 \\
\texttt{asyoulik.txt} & 122 & 4.81 & 4.86 & \textbf{3.51} & 3.12 & 228 & 107 \\
\texttt{bible.txt} & 3953 & 4.34 & 4.39 & \textbf{2.78} & 2.36 & 234 & 98 \\
\texttt{cp.html} & 24 & 5.23 & 5.35 & \textbf{3.68} & 2.59 & 231 & 89 \\
\texttt{fields.c} & 11 & 5.01 & 5.23 & \textbf{3.56} & 2.24 & 206 & 69 \\
\texttt{grammar.lsp} & 4 & 4.63 & 5.23 & \textbf{3.91} & 2.62 & 214 & 32 \\
\texttt{kennedy.xls} & 1006 & 3.57 & 3.60 & \textbf{2.40} & 1.58 & 286 & 142 \\
\texttt{lcet10.txt} & 417 & 4.67 & 4.70 & \textbf{3.19} & 2.72 & 235 & 102 \\
\texttt{plrabn12.txt} & 471 & 4.53 & 4.58 & \textbf{3.41} & 3.24 & 228 & 104 \\
\texttt{ptt5} & 501 & 1.21 & 1.66 & \textbf{0.97} & 0.88 & 251 & 203 \\
\texttt{sum} & 37 & 5.33 & 5.42 & \textbf{4.21} & 2.72 & 186 & 91 \\
\texttt{world192.txt} & 2415 & 5.00 & 5.04 & \textbf{3.26} & 2.34 & 196 & 96 \\
\texttt{xargs.1} & 4 & 4.90 & 5.42 & \textbf{4.43} & 3.28 & 91 & 5 \\
\midrule
\texttt{dickens} & 9954 & 4.53 & 4.57 & \textbf{3.29} & 3.04 & 216 & 95 \\
\texttt{mozilla} & 50020 & 6.22 & 6.24 & \textbf{4.30} & 2.98 & 150 & 63 \\
\texttt{mr} & 9737 & 3.68 & 3.71 & \textbf{3.01} & 2.95 & 246 & 112 \\
\texttt{nci} & 32767 & 2.43 & 2.44 & \textbf{0.94} & 0.76 & 264 & 137 \\
\texttt{ooffice} & 6008 & 6.64 & 6.66 & \textbf{5.29} & 4.03 & 205 & 83 \\
\texttt{osdb} & 9849 & 6.59 & 6.62 & \textbf{4.38} & 2.93 & 219 & 84 \\
\texttt{reymont} & 6472 & 4.84 & 4.87 & \textbf{2.46} & 2.25 & 218 & 106 \\
\texttt{samba} & 21100 & 6.09 & 6.13 & \textbf{3.06} & 2.02 & 212 & 100 \\
\texttt{sao} & 7082 & 7.53 & 7.55 & \textbf{7.53} & 5.88 & 214 & 80 \\
\texttt{webster} & 40487 & 4.97 & 5.00 & \textbf{2.76} & 2.36 & 207 & 95 \\
\texttt{x-ray} & 8276 & 6.60 & \textbf{6.63} & 6.67 & 5.71 & 237 & 102 \\
\texttt{xml} & 5220 & 5.52 & 5.55 & \textbf{1.66} & 1.03 & 247 & 103 \\

    \bottomrule
  \end{tabular}
  \par\smallskip
  \begin{minipage}{0.92\linewidth}
    \footnotesize $^\dagger$\,DEFLATE (zlib del JDK, nivel 6, LZ77 + Huffman) es solo una
    referencia externa y no forma parte de la comparación. Las columnas Huffman, LZW y DEFLATE
    están en bits por byte, incluida la cabecera.
  \end{minipage}
\end{table}}

En conjunto, sobre los \resCorpusFiles{} archivos, la tasa agregada es de
\resCorpusHuffBpb{}~bpb para Huffman, \resCorpusLzwBpb{}~bpb para LZW y
\resCorpusDeflateBpb{}~bpb para DEFLATE; LZW comprime mejor que Huffman en \resLzwWins{} de
los \resCorpusFiles{} archivos.

% TODO-RESULTS: interpretar tab:corpus (en qué tipos de archivo gana cada algoritmo, relación con H0, archivos donde LZW expande o se acerca a Huffman, rendimiento relativo en MB/s, distancia a DEFLATE).
\todoresults{Interpretar la \cref{tab:corpus}: tipos de archivo en que gana cada algoritmo,
relación con $H_0$, casos en que LZW se acerca o supera a 8~bpb, rendimiento relativo en MB/s
y distancia de ambos a DEFLATE.}
\end{aporte}

\subsection{E4: tamaño máximo del diccionario de LZW}
\label{sec:res-e4}
\begin{aporte}{C}
La \cref{fig:e4} muestra los bits por byte y el tiempo de codificación de LZW al variar el
ancho máximo de código $b$, es decir, el tamaño del diccionario $2^b$, en tres archivos de
naturaleza distinta. El ancho con menor bpb fue $b = \resLzwBitsBestDickens$ en
\texttt{dickens}, $b = \resLzwBitsBestXml$ en \texttt{xml} y $b = \resLzwBitsBestEcoli$ en
\texttt{E.coli}.

\siresultado{e4_lzw_bits.csv}{%
\begin{figure}[H]
  \centering
  \begin{subfigure}[t]{0.49\linewidth}
    \centering
    \begin{tikzpicture}
      \begin{axis}[informe, height=0.78\linewidth, xlabel={$b$ (bits; diccionario de $2^b$ entradas)},
          ylabel={bits por byte}, xtick={9,...,16}, legend pos=north east]
        \addplot[blue!70!black, mark=*, mark size=1.3pt]
          table[col sep=comma, x=bits, y=dickens_bpb] {results/e4_lzw_bits.csv};
        \addlegendentry{\texttt{dickens}}
        \addplot[green!45!black, mark=square*, mark size=1.2pt]
          table[col sep=comma, x=bits, y=xml_bpb] {results/e4_lzw_bits.csv};
        \addlegendentry{\texttt{xml}}
        \addplot[red!70!black, mark=triangle*, mark size=1.5pt]
          table[col sep=comma, x=bits, y=ecoli_bpb] {results/e4_lzw_bits.csv};
        \addlegendentry{\texttt{E.coli}}
      \end{axis}
    \end{tikzpicture}
    \caption{Tasa de compresión.}
    \label{fig:e4-bpb}
  \end{subfigure}\hfill
  \begin{subfigure}[t]{0.49\linewidth}
    \centering
    \begin{tikzpicture}
      \begin{axis}[informe, height=0.78\linewidth, xlabel={$b$ (bits; diccionario de $2^b$ entradas)},
          ylabel={Mediana de codificación [ms]}, xtick={9,...,16}, ymin=0, legend pos=north east]
        \addplot[blue!70!black, mark=*, mark size=1.3pt]
          table[col sep=comma, x=bits, y=dickens_enc_ms] {results/e4_lzw_bits.csv};
        \addlegendentry{\texttt{dickens}}
        \addplot[green!45!black, mark=square*, mark size=1.2pt]
          table[col sep=comma, x=bits, y=xml_enc_ms] {results/e4_lzw_bits.csv};
        \addlegendentry{\texttt{xml}}
        \addplot[red!70!black, mark=triangle*, mark size=1.5pt]
          table[col sep=comma, x=bits, y=ecoli_enc_ms] {results/e4_lzw_bits.csv};
        \addlegendentry{\texttt{E.coli}}
      \end{axis}
    \end{tikzpicture}
    \caption{Tiempo de codificación.}
    \label{fig:e4-tiempo}
  \end{subfigure}
  \caption{E4: efecto del tamaño máximo del diccionario de LZW.}
  \label{fig:e4}
\end{figure}}

% TODO-RESULTS: interpretar fig:e4 (tendencia del bpb con b, rendimientos decrecientes, diferencias entre texto, XML y ADN, costo temporal de diccionarios grandes y frecuencia de reinicios).
\todoresults{Interpretar la \cref{fig:e4}: tendencia de los bits por byte al crecer $b$,
rendimientos decrecientes, diferencias entre texto, XML y ADN, y costo temporal asociado al
tamaño del diccionario y a la frecuencia de reinicios.}
\end{aporte}
